\chapter*{摘要}
\addcontentsline{toc}{chapter}{摘要}

本文面向第二十一届全国大学生智能汽车竞赛飞跃雷区组，围绕江南大学博世·飞驰人生队的
空地协同系统，建立一份可持续补充、可追溯到仓库实现的技术报告框架。现有无人机端以
Infineon CYT4BB7 双核微控制器为核心：CM7\_0 负责飞行控制、状态估计、任务管理与安全逻辑，
CM7\_1 负责摄像头采集及视觉链路。系统已形成由视觉/红外定位、光流与 ToF 高度观测、
位置—速度—姿态—角速度多级闭环、DShot600 电机控制、空地通信和 VOFA 在线诊断构成的完整软件链路。

本报告后续将结合机械与电路实测资料，给出系统设计、传感器标定、控制算法、视觉定位、
空地协同、调试过程及性能实验。所有尚未由实测数据支撑的参数均在框架中显式标记，定稿时应以
真实测量结果、代码版本和可复现实验替换。

\todoitem{定稿时重写摘要：补充整机质量、定位误差、控制频率、任务成功率、平均完成时间等量化结果；摘要一般控制在一页内。}

\noindent\textbf{关键词：}飞跃雷区；空地协同；CYT4BB7；双核飞控；视觉定位；级联控制；光流；ToF
